% chapters/ch04.tex —— 第 4 章：最小二乘问题
% 源书：Chapter 4 Least Squares Problems，书页 46–59（物理页 60–73）
% 公式编号 (4.1)–(4.8)；定理 4.1–4.9；定义 4.1–4.3；推论 4.1–4.2；算法 4.1–4.2；习题 1–15
\chapter{最小二乘问题}
\label{ch:leastsquares}

本章研究线性最小二乘问题:
\[
  \min_{\mathbf{y}\in\R^n}\norm{A\mathbf{y}-\mathbf{b}}_2,
\]
其中数据矩阵 $A\in\R^{m\times n}$ ($m\geqslant n$) 与观测向量 $\mathbf{b}\in\R^m$ 是给定的. 为构造求解此类问题的高效算法, 我们将介绍一些著名的正交变换与 QR 分解. 关于最小二乘问题的文献, 可参见 \cite{src:9,src:11,src:28,src:36,src:88,src:92,src:93,src:94}.

\section{最小二乘问题}
\label{sec:4-1}

实践中, 若给定 $m$ 个点 $t_1,t_2,\ldots,t_m$ 以及这些点上的数据 $y_1,y_2,\ldots,y_m$, 并在这些点上定义函数 $\phi_1(t),\phi_2(t),\ldots,\phi_n(t)$, 我们便试图找出由
\[
  f(\mathbf{x},t)\equiv\sum_{j=1}^{n}x_j\phi_j(t)
\]
定义的函数 $f(\mathbf{x},t)$, 其中 $\mathbf{x}=[x_1,\ldots,x_n]^T$, 使得由
\[
  r_i(\mathbf{x})\equiv y_i-f(\mathbf{x},t_i)
  =y_i-\sum_{j=1}^{n}x_j\phi_j(t_i),
  \qquad i=1,2,\ldots,m
\]
定义的\newterm{residual}{残差}{residual}尽可能小. 写成矩阵形式, 有
\[
  \mathbf{r}(\mathbf{x})=\mathbf{b}-A\mathbf{x},
\]
其中
\[
  A=
  \begin{bmatrix}
    \phi_1(t_1) & \cdots & \phi_n(t_1)\\
    \vdots      &        & \vdots\\
    \phi_1(t_m) & \cdots & \phi_n(t_m)
  \end{bmatrix},
\]
以及
\[
  \mathbf{b}=[y_1,\ldots,y_m]^T,\qquad
  \mathbf{r}(\mathbf{x})=[r_1(\mathbf{x}),\ldots,r_m(\mathbf{x})]^T.
\]

当 $m=n$ 时, 我们可以要求 $\mathbf{r}(\mathbf{x})=\mathbf{0}$, 此时 $\mathbf{x}$ 可通过求解方程组 $A\mathbf{x}=\mathbf{b}$ 得到. 当 $m>n$ 时, 我们要求 $\mathbf{r}(\mathbf{x})$ 在范数 $\norm{\cdot}_2$ 下达到最小. 为此引入如下最小二乘问题的定义.

\begin{definition}
\label{def:4-1}
设 $A\in\R^{m\times n}$, $\mathbf{b}\in\R^m$. 求 $\mathbf{x}\in\R^n$ 使得
\begin{equation}
  \norm{\mathbf{b}-A\mathbf{x}}_2=\norm{\mathbf{r}(\mathbf{x})}_2
  =\min_{\mathbf{y}\in\R^n}\norm{\mathbf{r}(\mathbf{y})}_2
  =\min_{\mathbf{y}\in\R^n}\norm{\mathbf{b}-A\mathbf{y}}_2.
  \label{eq:4-1}
\end{equation}
该问题称为\newterm{leastsquares}{最小二乘问题}{least squares problem}, $\mathbf{r}(\mathbf{x})$ 称为残差.
\end{definition}

以下只考虑
\[
  \rank(A)=n<m
\]
的情形. 先研究如下方程的解 $\mathbf{x}$:
\begin{equation}
  A\mathbf{x}=\mathbf{b},\qquad A\in\R^{m\times n}.
  \label{eq:4-2}
\end{equation}

矩阵 $A$ 的\newterm{range}{值域}{range}定义为
\[
  \mathcal{R}(A)\equiv\{\mathbf{y}\in\R^m:\mathbf{y}=A\mathbf{x},\ \mathbf{x}\in\R^n\}.
\]
容易看出
\[
  \mathcal{R}(A)=\spn\{\mathbf{a}_1,\ldots,\mathbf{a}_n\},
\]
其中 $\mathbf{a}_i\ (i=1,\ldots,n)$ 是 $A$ 的列向量. $A$ 的\newterm{nullspace}{零空间}{nullspace}定义为
\[
  \mathcal{N}(A)\equiv\{\mathbf{x}\in\R^n:A\mathbf{x}=\mathbf{0}\}.
\]
$\mathcal{N}(A)$ 的维数称为 $A$ 的\newterm{nullity}{零度}{nullity}, 记为 $\operatorname{nullity}(A)$. 子空间 $\mathcal{S}\subset\R^n$ 的\newterm{orthcomplement}{正交补}{orthogonal complement}定义为
\[
  \mathcal{S}^{\perp}\equiv\{\mathbf{y}\in\R^n:
  \mathbf{y}^T\mathbf{x}=0,\ \text{对一切}\ \mathbf{x}\in\mathcal{S}\}.
\]

关于 \eqref{eq:4-2}, 我们有如下定理.

\begin{theorem}
\label{thm:4-1}
方程 \eqref{eq:4-2} 有解 $\iff\rank(A)=\rank([A,\mathbf{b}])$.
\end{theorem}

\begin{theorem}
\label{thm:4-2}
设 $\mathbf{x}$ 是 \eqref{eq:4-2} 的一个特解, 则 \eqref{eq:4-2} 的解集为 $\mathbf{x}+\mathcal{N}(A)$.
\end{theorem}

\begin{corollary}
\label{cor:4-1}
假设方程 \eqref{eq:4-2} 有解, 则解唯一 $\iff$ $\operatorname{nullity}(A)=0$.
\end{corollary}

关于 \eqref{eq:4-1} 的解, 我们有如下基本定理.

\begin{theorem}
\label{thm:4-3}
最小二乘问题 \eqref{eq:4-1} 总有解, 且解唯一当且仅当 $\operatorname{nullity}(A)=0$.
\end{theorem}

\begin{proof}
由于
\[
  \R^m=\mathcal{R}(A)\oplus\mathcal{R}(A)^{\perp},
\]
向量 $\mathbf{b}$ 可唯一地表示为
\[
  \mathbf{b}=\mathbf{b}_1+\mathbf{b}_2,
\]
其中 $\mathbf{b}_1\in\mathcal{R}(A)$, $\mathbf{b}_2\in\mathcal{R}(A)^{\perp}$. 对任意 $\mathbf{x}\in\R^n$, 由于 $\mathbf{b}_1-A\mathbf{x}\in\mathcal{R}(A)$ 且与 $\mathbf{b}_2$ 正交, 故有
\[
  \begin{aligned}
    \norm{\mathbf{r}(\mathbf{x})}_2^2
    &=\norm{\mathbf{b}-A\mathbf{x}}_2^2
     =\norm{(\mathbf{b}_1-A\mathbf{x})+\mathbf{b}_2}_2^2\\
    &=\norm{\mathbf{b}_1-A\mathbf{x}}_2^2+\norm{\mathbf{b}_2}_2^2.
  \end{aligned}
\]
注意 $\norm{\mathbf{r}(\mathbf{x})}_2^2$ 达到最小当且仅当 $\norm{\mathbf{b}_1-A\mathbf{x}}_2^2$ 达到最小. 由于 $\mathbf{b}_1\in\mathcal{R}(A)$, $\norm{\mathbf{r}(\mathbf{x})}_2^2$ 达到最小当且仅当
\[
  A\mathbf{x}=\mathbf{b}_1,
\]
即
\[
  \norm{\mathbf{b}_1-A\mathbf{x}}_2^2=0.
\]
于是由推论 \ref{cor:4-1} 知, $A\mathbf{x}=\mathbf{b}_1$ 的解唯一 (即 \eqref{eq:4-1} 的解唯一) 当且仅当
\[
  \operatorname{nullity}(A)=0. \qedhere
\]
\end{proof}

令
\[
  \mathcal{X}=\{\mathbf{x}\in\R^n:\mathbf{x}\ \text{是 \eqref{eq:4-1} 的解}\}.
\]
我们有

\begin{theorem}
\label{thm:4-4}
向量 $\mathbf{x}\in\mathcal{X}$ 当且仅当
\begin{equation}
  A^TA\mathbf{x}=A^T\mathbf{b}.
  \label{eq:4-3}
\end{equation}
\end{theorem}

\begin{proof}
设 $\mathbf{x}\in\mathcal{X}$. 由定理 \ref{thm:4-3} 知 $A\mathbf{x}=\mathbf{b}_1$, 其中 $\mathbf{b}_1\in\mathcal{R}(A)$, 且
\[
  \mathbf{r}(\mathbf{x})=\mathbf{b}-A\mathbf{x}=\mathbf{b}-\mathbf{b}_1
  =\mathbf{b}_2\in\mathcal{R}(A)^{\perp}.
\]
因此
\[
  A^T\mathbf{r}(\mathbf{x})=A^T\mathbf{b}_2=\mathbf{0}.
\]
把 $\mathbf{r}(\mathbf{x})=\mathbf{b}-A\mathbf{x}$ 代入 $A^T\mathbf{r}(\mathbf{x})=\mathbf{0}$, 即得 \eqref{eq:4-3}.

反之, 设 $\mathbf{x}\in\R^n$ 满足
\[
  A^TA\mathbf{x}=A^T\mathbf{b},
\]
则对任意 $\mathbf{y}\in\R^n$, 有
\[
  \begin{aligned}
    \norm{\mathbf{b}-A(\mathbf{x}+\mathbf{y})}_2^2
    &=\norm{\mathbf{b}-A\mathbf{x}}_2^2
     -2\mathbf{y}^TA^T(\mathbf{b}-A\mathbf{x})+\norm{A\mathbf{y}}_2^2\\
    &=\norm{\mathbf{b}-A\mathbf{x}}_2^2+\norm{A\mathbf{y}}_2^2\\
    &\geqslant\norm{\mathbf{b}-A\mathbf{x}}_2^2.
  \end{aligned}
\]
于是 $\mathbf{x}\in\mathcal{X}$. \qedhere
\end{proof}

因此, 求解最小二乘问题有如下算法:
\begin{parenlist}
  \item 计算 $C=A^TA$ 与 $\mathbf{d}=A^T\mathbf{b}$.
  \item 求 $C=LL^T$ 的 Cholesky 分解.
  \item 求解三角线性方程组: $L\mathbf{y}=\mathbf{d}$ 与 $L^T\mathbf{x}=\mathbf{y}$.
\end{parenlist}

我们指出, 计算 $A^TA$ 的运算量通常为 $\mathcal{O}(n^2m)$, 并且 $A$ 的某些信息可能丢失. 例如, 考虑
\[
  A=
  \begin{bmatrix}
    1 & 1 & 1\\
    \epsilon & 0 & 0\\
    0 & \epsilon & 0\\
    0 & 0 & \epsilon
  \end{bmatrix}.
\]
有
\[
  A^TA=
  \begin{bmatrix}
    1+\epsilon^2 & 1 & 1\\
    1 & 1+\epsilon^2 & 1\\
    1 & 1 & 1+\epsilon^2
  \end{bmatrix}.
\]
假设 $\epsilon=10^{-3}$ 且采用 6 位十进制浮点系统, 则 $1+\epsilon^2=1+10^{-6}$ 被舍入为 $1$, 这意味着 $A^TA$ 是奇异的!

我们指出, \eqref{eq:4-3} 的解 $\mathbf{x}$ 可表示为
\[
  \mathbf{x}=(A^TA)^{-1}A^T\mathbf{b}.
\]
若令
\[
  A^{\dagger}=(A^TA)^{-1}A^T,
\]
则最小二乘解 $\mathbf{x}$ 可写为
\[
  \mathbf{x}=A^{\dagger}\mathbf{b}.
\]
实际上, $n\times m$ 矩阵 $A^{\dagger}$ 就是 $A$ 的 Moore–Penrose 广义逆, 它是唯一的 \cite{src:30,src:39,src:88}. 一般而言, 有

\begin{definition}
\label{def:4-2}
设 $X\in\R^{n\times m}$. 若它满足以下条件:
\[
  AXA=A,\qquad XAX=X,\qquad (AX)^T=AX,\qquad (XA)^T=XA,
\]
则称 $X$ 为 $A$ 的\newterm{moorepenrose}{Moore–Penrose 广义逆}{Moore–Penrose generalized inverse}, 记为 $A^{\dagger}$.
\end{definition}

下面建立最小二乘问题的扰动分析. 假设 $\mathbf{b}$ 上有一个扰动 $\delta\mathbf{b}$, 并记 $\mathbf{x}$ 与 $\mathbf{x}+\delta\mathbf{x}$ 分别为如下最小二乘问题的解:
\[
  \min_{\mathbf{x}}\norm{\mathbf{b}-A\mathbf{x}}_2,
  \qquad
  \min_{\mathbf{x}}\norm{(\mathbf{b}+\delta\mathbf{b})-A\mathbf{x}}_2.
\]
则
\[
  \mathbf{x}=A^{\dagger}\mathbf{b},
\]
且
\[
  \mathbf{x}+\delta\mathbf{x}
  =A^{\dagger}(\mathbf{b}+\delta\mathbf{b})=A^{\dagger}\tilde{\mathbf{b}},
\]
其中 $\tilde{\mathbf{b}}=\mathbf{b}+\delta\mathbf{b}$. 矩阵 $A\in\R^{m\times n}$ 的矩阵 2 范数定义为
\[
  \norm{A}_2\equiv\max_{\mathbf{0}\neq\mathbf{x}\in\R^n}
  \frac{\norm{A\mathbf{x}}_2}{\norm{\mathbf{x}}_2}
  =\max_{\norm{\mathbf{x}}_2=1}\norm{A\mathbf{x}}_2,
\]
这与方阵的矩阵 2 范数相同. 此处有

\begin{theorem}
\label{thm:4-5}
设 $\mathbf{b}_1$ 与 $\tilde{\mathbf{b}}_1$ 分别表示 $\mathbf{b}$ 与 $\tilde{\mathbf{b}}$ 在 $\mathcal{R}(A)$ 上的\newterm{orthprojection}{正交投影}{orthogonal projection}. 若 $\mathbf{b}_1\neq\mathbf{0}$, 则
\[
  \frac{\norm{\delta\mathbf{x}}_2}{\norm{\mathbf{x}}_2}
  \leqslant\kappa_2(A)\frac{\norm{\delta\mathbf{b}_1}_2}{\norm{\mathbf{b}_1}_2},
\]
其中 $\kappa_2(A)=\norm{A}_2\norm{A^{\dagger}}_2$, 且 $\tilde{\mathbf{b}}_1=\mathbf{b}_1+\delta\mathbf{b}_1$.
\end{theorem}

\begin{proof}
设 $\mathbf{b}_2$ 表示 $\mathbf{b}$ 在 $\mathcal{R}(A)^{\perp}$ 上的正交投影, 则 $\mathbf{b}=\mathbf{b}_1+\mathbf{b}_2$ 且 $A^T\mathbf{b}_2=\mathbf{0}$. 注意
\[
  A^{\dagger}\mathbf{b}=A^{\dagger}\mathbf{b}_1+A^{\dagger}\mathbf{b}_2
  =A^{\dagger}\mathbf{b}_1+(A^TA)^{-1}A^T\mathbf{b}_2
  =A^{\dagger}\mathbf{b}_1.
\]
类似地, $A^{\dagger}\tilde{\mathbf{b}}=A^{\dagger}\tilde{\mathbf{b}}_1$. 因此
\begin{align}
  \norm{\delta\mathbf{x}}_2
  &=\norm{A^{\dagger}\mathbf{b}-A^{\dagger}\tilde{\mathbf{b}}}_2
   =\norm{A^{\dagger}(\mathbf{b}_1-\tilde{\mathbf{b}}_1)}_2\notag\\
  &\leqslant\norm{A^{\dagger}}_2\norm{\mathbf{b}_1-\tilde{\mathbf{b}}_1}_2
   =\norm{A^{\dagger}}_2\norm{\delta\mathbf{b}_1}_2.
   \label{eq:4-4}
\end{align}
由于 $A\mathbf{x}=\mathbf{b}_1$, 有
\begin{equation}
  \norm{\mathbf{b}_1}_2\leqslant\norm{A}_2\norm{\mathbf{x}}_2.
  \label{eq:4-5}
\end{equation}
结合 \eqref{eq:4-4} 与 \eqref{eq:4-5}, 证毕.
\end{proof}

我们指出, 条件数 $\kappa_2(A)$ 对最小二乘问题十分重要. 当 $\kappa_2(A)$ 大时, 称该最小二乘问题是病态的; 当 $\kappa_2(A)$ 小时, 称该最小二乘问题是良态的.

\begin{theorem}
\label{thm:4-6}
设 $A$ 的列向量线性无关, 则
\[
  \kappa_2(A)^2=\kappa_2(A^TA).
\]
\end{theorem}

\begin{proof}
由定理 \ref{thm:3-3} (ii) 与所给条件, 有
\[
  \norm{A}_2^2=\norm{A^TA}_2,\qquad
  \norm{A^{\dagger}}_2^2=\norm{A^{\dagger}(A^{\dagger})^T}_2
  =\norm{(A^TA)^{-1}}_2.
\]
因此
\[
  \kappa_2(A)^2=\norm{A}_2^2\norm{A^{\dagger}}_2^2
  =\norm{A^TA}_2\norm{(A^TA)^{-1}}_2=\kappa_2(A^TA). \qedhere
\]
\end{proof}

\section{正交变换}
\label{sec:4-2}

为构造求解最小二乘问题的高效算法, 我们介绍一些著名的正交变换.

\subsection{Householder 变换}
\label{subsec:4-2-1}

首先介绍 Householder 变换的如下定义.

\begin{definition}
\label{def:4-3}
设 $\mathbf{w}\in\R^n$ 且 $\norm{\mathbf{w}}_2=1$. 定义 $H\in\R^{n\times n}$ 如下:
\begin{equation}
  H=I-2\mathbf{w}\mathbf{w}^T.
  \label{eq:4-6}
\end{equation}
矩阵 $H$ 称为 \newterm{householder}{Householder 变换}{Householder transformation}.
\end{definition}

\begin{theorem}
\label{thm:4-7}
设 $H$ 由 \eqref{eq:4-6} 定义, 则 $H$ 具有以下性质:
\begin{romanlist}
  \item $H$ 是对称正交矩阵.
  \item $H^2=I$.
  \item $H$ 称为\newterm{reflection}{反射}{reflection}, 因为 $H\mathbf{x}$ 是 $\mathbf{x}\in\R^n$ 关于过 $\mathbf{0}$ 且垂直于 $\mathbf{w}$ 的平面的反射.
\end{romanlist}
\end{theorem}

\begin{proof}
我们只证 (iii). 注意对任意向量 $\mathbf{x}\in\R^n$, 它可表示为
\[
  \mathbf{x}=\mathbf{u}+\alpha\mathbf{w},
\]
其中 $\mathbf{u}\in\spn\{\mathbf{w}\}^{\perp}$, $\alpha\in\R$. 利用 $\mathbf{u}^T\mathbf{w}=0$ 与 $\mathbf{w}^T\mathbf{w}=1$, 有
\[
  H\mathbf{x}=(I-2\mathbf{w}\mathbf{w}^T)(\mathbf{u}+\alpha\mathbf{w})
  =\mathbf{u}+\alpha\mathbf{w}-2\mathbf{w}\mathbf{w}^T\mathbf{u}
  -2\alpha\mathbf{w}\mathbf{w}^T\mathbf{w}
  =\mathbf{u}-\alpha\mathbf{w}. \qedhere
\]
\end{proof}

\begin{theorem}
\label{thm:4-8}
对任意 $\mathbf{0}\neq\mathbf{x}\in\R^n$, 可以构造单位向量 $\mathbf{w}$, 使得由 \eqref{eq:4-6} 定义的 Householder 变换满足
\[
  H\mathbf{x}=\alpha\mathbf{e}_1,
\]
其中 $\alpha=\pm\norm{\mathbf{x}}_2$.
\end{theorem}

\begin{proof}
注意 $H\mathbf{x}=(I-2\mathbf{w}\mathbf{w}^T)\mathbf{x}
=\mathbf{x}-2(\mathbf{w}^T\mathbf{x})\mathbf{w}$. 令
\[
  \mathbf{w}=\frac{\mathbf{x}-\alpha\mathbf{e}_1}
  {\norm{\mathbf{x}-\alpha\mathbf{e}_1}_2},
\]
则有
\begin{align}
  H\mathbf{x}&=\mathbf{x}-2(\mathbf{w}^T\mathbf{x})\mathbf{w}\notag\\
  &=\mathbf{x}-\frac{2}{\norm{\mathbf{x}-\alpha\mathbf{e}_1}_2^2}
   \bigl[(\mathbf{x}^T-\alpha\mathbf{e}_1^T)\mathbf{x}\bigr]
   (\mathbf{x}-\alpha\mathbf{e}_1)\notag\\
  &=\mathbf{x}
   -\frac{2\bigl(\norm{\mathbf{x}}_2^2-\alpha x_1\bigr)}
   {\norm{\mathbf{x}-\alpha\mathbf{e}_1}_2^2}\mathbf{x}
   +\frac{2\bigl(\norm{\mathbf{x}}_2^2-\alpha x_1\bigr)\alpha}
   {\norm{\mathbf{x}-\alpha\mathbf{e}_1}_2^2}\mathbf{e}_1\notag\\
  &=\left[1-\frac{2\bigl(\norm{\mathbf{x}}_2^2-\alpha x_1\bigr)}
   {\norm{\mathbf{x}-\alpha\mathbf{e}_1}_2^2}\right]\mathbf{x}
   +\frac{2\bigl(\norm{\mathbf{x}}_2^2-\alpha x_1\bigr)\alpha}
   {\norm{\mathbf{x}-\alpha\mathbf{e}_1}_2^2}\mathbf{e}_1,
   \label{eq:4-7}
\end{align}
其中 $x_1$ 是向量 $\mathbf{x}$ 的第一个分量. 令 $\mathbf{x}$ 的系数为零, 便得如下方程:
\[
  1-\frac{2\bigl(\norm{\mathbf{x}}_2^2-\alpha x_1\bigr)}
  {\norm{\mathbf{x}-\alpha\mathbf{e}_1}_2^2}=0.
\]
解此方程得 $\alpha=\pm\norm{\mathbf{x}}_2$. 把它代入 \eqref{eq:4-7}, 便有
\[
  H\mathbf{x}=\mp\norm{\mathbf{x}}_2\mathbf{e}_1. \qedhere
\]
\end{proof}

我们指出, 对任意向量 $\mathbf{0}\neq\mathbf{x}\in\R^n$, 由定理 \ref{thm:4-8}, 可以构造 Householder 矩阵 $H$, 使 $H\mathbf{x}$ 的后 $n-1$ 个分量均为零. 可用以下两步构造 $H$ 的单位向量 $\mathbf{w}$:
\begin{parenlist}
  \item 计算 $\mathbf{v}=\mathbf{x}\pm\norm{\mathbf{x}}_2\mathbf{e}_1$;
  \item 计算 $\mathbf{w}=\mathbf{v}/\norm{\mathbf{v}}_2$.
\end{parenlist}

接下来的自然问题是: $\norm{\mathbf{x}}_2$ 前面的符号如何选取? 通常取
\[
  \mathbf{v}=\mathbf{x}+\sign(x_1)\norm{\mathbf{x}}_2\mathbf{e}_1,
\]
其中 $x_1\neq0$ 是向量 $\mathbf{x}$ 的第一个分量 \cite{src:80}. 由于
\[
  H=I-2\mathbf{w}\mathbf{w}^T
  =I-\frac{2}{\mathbf{v}^T\mathbf{v}}\mathbf{v}\mathbf{v}^T
  =I-\beta\mathbf{v}\mathbf{v}^T,
\]
其中 $\beta=2/\mathbf{v}^T\mathbf{v}$, 我们只需计算 $\beta$ 与 $\mathbf{v}$, 而不必形成 $\mathbf{w}$. 于是有如下算法.

\begin{algorithm}{Householder 变换}
\label{alg:4-1}
\[
  \left\{
  \begin{array}{l}
    \text{\textbf{function}:}\ [\mathbf{v},\beta]=\text{house}(\mathbf{x})\\
    \quad n=\text{length}(\mathbf{x})\\
    \quad\sigma=\mathbf{x}(2:n)^T\mathbf{x}(2:n)\\
    \quad\mathbf{v}(1)=\mathbf{x}(1)+\sign(\mathbf{x}(1))
      \sqrt{\mathbf{x}(1)^2+\sigma}\\
    \quad\mathbf{v}(2:n)=\mathbf{x}(2:n)\\
    \quad\textbf{if}\ \sigma=0\\
    \quad\quad\beta=0\\
    \quad\textbf{else}\\
    \quad\quad\beta=2/(\mathbf{v}(1)^2+\sigma)\\
    \quad\textbf{end}
  \end{array}
  \right.
\]
\end{algorithm}

\subsection{Givens 旋转}
\label{subsec:4-2-2}

\newterm{givens}{Givens 旋转}{Givens rotation}定义如下:
\[
  \begin{aligned}
    G(i,k,\theta)
    &=I+s(\mathbf{e}_i\mathbf{e}_k^T-\mathbf{e}_k\mathbf{e}_i^T)
      +(c-1)(\mathbf{e}_i\mathbf{e}_i^T+\mathbf{e}_k\mathbf{e}_k^T)\\
    &=
    \begin{bmatrix}
      1      &        & \vdots &        & \vdots &        &        \\
             & \ddots & \vdots &        & \vdots &        &        \\
      \cdots & \cdots & c      & \cdots & s      & \cdots & \cdots \\
             &        & \vdots &        & \vdots &        &        \\
      \cdots & \cdots & -s     & \cdots & c      & \cdots & \cdots \\
             &        & \vdots &        & \vdots & \ddots &        \\
             &        & \vdots &        & \vdots &        & 1
    \end{bmatrix},
  \end{aligned}
\]
其中 $c=\cos\theta$, $s=\sin\theta$. 容易证明 $G(i,k,\theta)$ 是正交矩阵.

设 $\mathbf{x}=[x_1,\ldots,x_n]^T\in\R^n$, $\mathbf{y}=G(i,k,\theta)\mathbf{x}$, 则有
\[
  y_i=cx_i+sx_k,\qquad y_k=-sx_i+cx_k,\qquad
  y_j=x_j,\quad j\neq i,k.
\]
若要使 $y_k=0$, 只需取
\[
  c=\frac{x_i}{\sqrt{x_i^2+x_k^2}},\qquad
  s=\frac{x_k}{\sqrt{x_i^2+x_k^2}}.
\]
因此
\[
  y_i=\sqrt{x_i^2+x_k^2},\qquad y_k=0.
\]
我们指出, 对任意向量 $\mathbf{0}\neq\mathbf{x}\in\R^n$, 可以构造作用在 $\mathbf{x}$ 上的 Givens 旋转 $G(i,k,\theta)$, 使 $\mathbf{x}$ 的某个非零分量变为零.

\section{QR 分解}
\label{sec:4-3}

设 $A\in\R^{m\times n}$, $\mathbf{b}\in\R^m$. 由定理 \ref{thm:3-3} (iii), 对任意正交矩阵 $Q$, 有
\[
  \norm{A\mathbf{x}-\mathbf{b}}_2=\norm{Q^T(A\mathbf{x}-\mathbf{b})}_2.
\]
因此, 最小二乘问题
\[
  \min_{\mathbf{x}}\norm{Q^TA\mathbf{x}-Q^T\mathbf{b}}_2
\]
与 \eqref{eq:4-1} 等价. 我们希望找到一个合适的正交矩阵 $Q$, 使原来的最小二乘问题化为更易求解的最小二乘问题. 我们有

\begin{theorem}
\label{thm:4-9}
(QR 分解) 设 $A\in\R^{m\times n}$ 且 $m\geqslant n$, 则 $A$ 具有 QR 分解:
\begin{equation}
  A=Q
  \begin{bmatrix}
    R\\
    \mathbf{0}
  \end{bmatrix},
  \label{eq:4-8}
\end{equation}
其中 $Q\in\R^{m\times m}$ 是正交矩阵, $R\in\R^{n\times n}$ 是对角元非负的上三角矩阵. 当 $m=n$ 且 $A$ 非奇异时, 该分解唯一.
\end{theorem}

\begin{proof}
对 $n$ 作归纳. 当 $n=1$ 时, 由定理 \ref{thm:4-8} 知结论成立. 今设定理对 $\R^{p\times(n-1)}$ ($p\geqslant n-1$) 中的所有矩阵成立. 设 $A\in\R^{m\times n}$ 的第一列为 $\mathbf{a}_1$. 再次由定理 \ref{thm:4-8}, 存在正交矩阵 $Q_1\in\R^{m\times m}$ 使得
\[
  Q_1^T\mathbf{a}_1=\norm{\mathbf{a}_1}_2\mathbf{e}_1.
\]
因此, 有
\[
  Q_1^TA=
  \begin{bmatrix}
    \norm{\mathbf{a}_1}_2 & \mathbf{v}^T\\
    \mathbf{0}            & A_1
  \end{bmatrix}.
\]
对矩阵 $A_1\in\R^{(m-1)\times(n-1)}$, 由归纳假设得
\[
  A_1=Q_2
  \begin{bmatrix}
    R_2\\
    \mathbf{0}
  \end{bmatrix},
\]
其中 $Q_2\in\R^{(m-1)\times(m-1)}$ 是正交矩阵, $R_2$ 是对角元非负的上三角矩阵. 于是令
\[
  Q=Q_1
  \begin{bmatrix}
    1 & \mathbf{0}\\
    \mathbf{0} & Q_2
  \end{bmatrix},
  \qquad
  R=
  \begin{bmatrix}
    \norm{\mathbf{a}_1}_2 & \mathbf{v}^T\\
    \mathbf{0}            & R_2\\
    \mathbf{0}            & \mathbf{0}
  \end{bmatrix}.
\]
则 $Q$ 与 $R$ 即为满足定理条件的矩阵.

当 $A\in\R^{m\times m}$ 非奇异时, 我们要证明 QR 分解唯一. 设
\[
  A=QR=\widetilde{Q}\widetilde{R},
\]
其中 $Q,\widetilde{Q}\in\R^{m\times m}$ 是正交矩阵, $R,\widetilde{R}\in\R^{m\times m}$ 是对角元非负的上三角矩阵. 由于 $A$ 非奇异, $R$ 与 $\widetilde{R}$ 的对角元均为正. 因此, 矩阵
\[
  \widetilde{Q}^TQ=\widetilde{R}R^{-1}
\]
既正交又是对角元为正的上三角矩阵. 于是
\[
  \widetilde{Q}^TQ=\widetilde{R}R^{-1}=I,
\]
即
\[
  \widetilde{Q}=Q,\qquad\widetilde{R}=R. \qedhere
\]
\end{proof}

后面需要 QR 分解的复矩阵版本.

\begin{corollary}
\label{cor:4-2}
设 $A\in\C^{m\times n}$ 且 $m\geqslant n$, 则 $A$ 具有 QR 分解:
\[
  A=Q
  \begin{bmatrix}
    R\\
    \mathbf{0}
  \end{bmatrix},
\]
其中 $Q\in\C^{m\times m}$ 是\newterm{unitarymatrix}{酉矩阵}{unitary matrix}, $R\in\C^{n\times n}$ 是对角元非负的上三角矩阵. 当 $m=n$ 且 $A$ 非奇异时, 该分解唯一.
\end{corollary}

下面用 QR 分解求解最小二乘问题 \eqref{eq:4-1}. 假设 $A\in\R^{m\times n}$ ($m\geqslant n$) 的列线性无关, $\mathbf{b}\in\R^m$, 且 $A$ 具有 QR 分解 \eqref{eq:4-8}. 把 $Q$ 分块为
\[
  Q=[Q_1\mid Q_2],
\]
以及
\[
  Q^T\mathbf{b}=
  \begin{bmatrix}
    Q_1^T\\
    Q_2^T
  \end{bmatrix}\mathbf{b}
  =
  \begin{bmatrix}
    \mathbf{c}_1\\
    \mathbf{c}_2
  \end{bmatrix}.
\]
则
\[
  \norm{A\mathbf{x}-\mathbf{b}}_2^2
  =\norm{Q^TA\mathbf{x}-Q^T\mathbf{b}}_2^2
  =\norm{R\mathbf{x}-\mathbf{c}_1}_2^2+\norm{\mathbf{c}_2}_2^2.
\]
$\mathbf{x}$ 是最小二乘问题 \eqref{eq:4-1} 的解当且仅当它是 $R\mathbf{x}=\mathbf{c}_1$ 的解. 注意, 由于 $R$ 是上三角矩阵, 通过求解 $R\mathbf{x}=\mathbf{c}_1$ 来得到 \eqref{eq:4-1} 的解要容易得多. 于是有如下求解最小二乘问题的算法:
\begin{parenlist}
  \item 计算 $A$ 的 QR 分解.
  \item 计算 $\mathbf{c}_1=Q_1^T\mathbf{b}$.
  \item 求解上三角方程组 $R\mathbf{x}=\mathbf{c}_1$.
\end{parenlist}

最后, 我们讨论如何用 Householder 变换计算 $A$ 的 QR 分解. 设 $m=7$, $n=5$. 假设已求出 Householder 变换 $H_1$ 与 $H_2$, 使得
\[
  H_2H_1A=
  \begin{bmatrix}
    \times & \times & \times & \times & \times\\
    0      & \times & \times & \times & \times\\
    0      & 0      & +      & \times & \times\\
    0      & 0      & +      & \times & \times\\
    0      & 0      & +      & \times & \times\\
    0      & 0      & +      & \times & \times\\
    0      & 0      & +      & \times & \times
  \end{bmatrix}.
\]
现在构造 Householder 变换 $\widetilde{H}_3\in\R^{5\times5}$, 使得
\[
  \widetilde{H}_3
  \begin{bmatrix}
    +\\ +\\ +\\ +\\ +
  \end{bmatrix}
  =
  \begin{bmatrix}
    \times\\ 0\\ 0\\ 0\\ 0
  \end{bmatrix}.
\]
令 $H_3=\diag(I_2,\widetilde{H}_3)$, 得到
\[
  H_3H_2H_1A=
  \begin{bmatrix}
    \times & \times & \times & \times & \times\\
    0      & \times & \times & \times & \times\\
    0      & 0      & \times & \times & \times\\
    0      & 0      & 0      & \times & \times\\
    0      & 0      & 0      & \times & \times\\
    0      & 0      & 0      & \times & \times\\
    0      & 0      & 0      & \times & \times
  \end{bmatrix}.
\]
一般地, 经过 $n$ 个这样的步骤之后, 可把矩阵 $A$ 化为如下形式:
\[
  H_nH_{n-1}\cdots H_1A=
  \begin{bmatrix}
    R\\
    \mathbf{0}
  \end{bmatrix},
\]
其中 $R$ 是对角元非负的上三角矩阵. 令 $Q=H_1\cdots H_n$, 即得
\[
  A=Q
  \begin{bmatrix}
    R\\
    \mathbf{0}
  \end{bmatrix}.
\]
于是有如下算法.

\begin{algorithm}{QR 分解: Householder 变换}
\label{alg:4-2}
\[
  \left\{
  \begin{array}{l}
    \textbf{for}\ j=1:n\\
    \quad[\mathbf{v},\beta]=\text{house}(A(j:m,j))\\
    \quad A(j:m,j:n)=(I_{m-j+1}-\beta\mathbf{v}\mathbf{v}^T)A(j:m,j:n)\\
    \quad\textbf{if}\ j<m\\
    \quad\quad A(j+1:m,j)=\mathbf{v}(2:m-j+1)\\
    \quad\textbf{end}\\
    \textbf{end}
  \end{array}
  \right.
\]
\end{algorithm}

我们指出, QR 分解不仅是求解最小二乘问题的基本工具, 也是求解数值线性代数中其他一些基本问题的重要工具 \cite{src:20,src:23,src:24}.

\begin{exercises}
  \item 设 $A\in\R^{m\times n}$ 列满秩. 证明: 若 $E$ 满足 $\norm{E}_2\leqslant1/\norm{A^{\dagger}}_2$, 则 $A+E$ 也列满秩, 其中 $A^{\dagger}=(A^TA)^{-1}A^T$.

  \item 设 $U=[u_{ij}]$ 为非奇异上三角矩阵. 证明
  \[
    \kappa_\infty(U)\geqslant
    \frac{\underset{i}{\max}|u_{ii}|}{\underset{i}{\min}|u_{ii}|},
  \]
  其中 $\kappa_\infty(U)=\norm{U}_\infty\norm{U^{-1}}_\infty$.

  \item 设 $A\in\R^{m\times n}$, $m\geqslant n$, 且列满秩. 证明
  \[
    \begin{bmatrix}
      I & A\\
      A^T & \mathbf{0}
    \end{bmatrix}
    \begin{bmatrix}
      \mathbf{r}\\
      \mathbf{x}
    \end{bmatrix}
    =
    \begin{bmatrix}
      \mathbf{b}\\
      \mathbf{0}
    \end{bmatrix}
  \]
  有解, 其中 $\mathbf{x}$ 使 $\norm{A\mathbf{x}-\mathbf{b}}_2$ 达到最小.

  \item 设 $\mathbf{x}\in\R^n$, $P$ 是 Householder 变换且满足
  \[
    P\mathbf{x}=\pm\norm{\mathbf{x}}_2\mathbf{e}_1.
  \]
  设 $G_{12},G_{23},\ldots,G_{n-1,n}$ 为 Givens 旋转, 且令
  \[
    Q=G_{12}G_{23}\cdots G_{n-1,n}.
  \]
  若 $Q\mathbf{x}=\pm\norm{\mathbf{x}}_2\mathbf{e}_1$, 问 $P$ 是否等于 $Q$? 给出证明或反例.

  \item 设 $A\in\R^{m\times n}$. 证明 $X=A^{\dagger}$ 在所有 $X\in\R^{n\times m}$ 上使 $\norm{AX-I}_F$ 达到最小. 最小值是多少?

  \item 设
  $\mathbf{x}=
  \begin{bmatrix}
    x_1\\ x_2
  \end{bmatrix}\in\C^2$.
  给出计算如下酉矩阵的算法
  \[
    Q=
    \begin{bmatrix}
      c & \bar{s}\\
      -s & c
    \end{bmatrix},
    \qquad c\in\R,\quad c^2+|s|^2=1,
  \]
  使得
  $Q\mathbf{x}=
  \begin{bmatrix}
    *\\ 0
  \end{bmatrix}$.

  \item 设 $m\times n$ 矩阵 $A$ 具有形式
  \[
    A=
    \begin{bmatrix}
      A_1\\
      A_2
    \end{bmatrix},
  \]
  其中 $A_1$ 是 $n\times n$ 非奇异矩阵, $A_2$ 是 $(m-n)\times n$ 任意矩阵. 证明 $\norm{A^{\dagger}}_2\leqslant\norm{A_1^{-1}}_2$.

  \item 考虑如下著名的病态矩阵
  \[
    A=
    \begin{bmatrix}
      1 & 1 & 1\\
      \epsilon & 0 & 0\\
      0 & \epsilon & 0\\
      0 & 0 & \epsilon
    \end{bmatrix},
    \qquad|\epsilon|\ll1.
  \]
  \par\smallskip
  \hangindent=2em\hangafter=1\noindent\hspace*{2em}(a)\enspace
  选取小的 $\epsilon$ 使 $\rank(A)=3$, 然后计算 $\kappa_2(A)$, 以说明 $A$ 是病态的.
  \par\smallskip
  \hangindent=2em\hangafter=1\noindent\hspace*{2em}(b)\enspace
  对如上给定的 $A$ 与 $\mathbf{b}=[3,\epsilon,\epsilon,\epsilon]^T$, 分别用以下方法求最小二乘解:
  \par\smallskip
  \hangindent=4em\hangafter=1\noindent\hspace*{4em}(i)\enspace 正规方程法;
  \par\smallskip
  \hangindent=4em\hangafter=1\noindent\hspace*{4em}(ii)\enspace QR 方法.

  \item 设 $A=BC$, 其中 $B\in\C^{m\times r}$, $C\in\C^{r\times n}$, 且
  \[
    r=\rank(A)=\rank(B)=\rank(C).
  \]
  证明
  \[
    A^{\dagger}=C^*(CC^*)^{-1}(B^*B)^{-1}B^*.
  \]

  \item 设 $A=U\Sigma V^*\in\C^{m\times n}$, 其中 $U\in\C^{m\times n}$ 满足 $U^*U=I$, $V\in\C^{n\times n}$ 满足 $V^*V=I$, $\Sigma$ 是 $n\times n$ 对角矩阵. 证明
  \[
    A^{\dagger}=V\Sigma^{\dagger}U^*.
  \]

  \item 证明
  \[
    A^{\dagger}=\lim_{\epsilon\to0}(A^*A+\epsilon I)^{-1}A^*
    =\lim_{\epsilon\to0}A^*(AA^*+\epsilon I)^{-1}.
  \]

  \item 证明
  \[
    \mathcal{R}(A)\cap\mathcal{N}(A^*)=\{\mathbf{0}\}.
  \]

  \item 设 $A=[a_{ij}]\in\C^{n\times n}$ 为\newterm{idempotent}{幂等的}{idempotent}矩阵, 则
  \[
    \mathcal{R}(A)\oplus\mathcal{N}(A)=\C^n,\qquad
    \rank(A)=\sum_{i=1}^{n}a_{ii}.
  \]

  \item 设 $A\in\C^{m\times n}$. 证明
  \[
    \mathcal{R}(AA^{\dagger})=\mathcal{R}(AA^*)=\mathcal{R}(A),
  \]
  \[
    \mathcal{R}(A^{\dagger}A)=\mathcal{R}(A^*A)
    =\mathcal{R}(A^{\dagger})=\mathcal{R}(A^*),
  \]
  \[
    \mathcal{N}(AA^{\dagger})=\mathcal{N}(AA^*)
    =\mathcal{N}(A^{\dagger})=\mathcal{N}(A^*),
  \]
  \[
    \mathcal{N}(A^{\dagger}A)=\mathcal{N}(A^*A)=\mathcal{N}(A).
  \]
  因此 $AA^{\dagger}$ 与 $A^{\dagger}A$ 是正交投影算子.

  \item 证明推论 \ref{cor:4-2}.
\end{exercises}
